\section{Implementation}
\label{sec:implementation}

The implementation follows the receiver's control flow. The kernel first
validates the capability, scope, revocation state, and guards. A
cross-organization request then enters
\codepath{ChioRuntimeAdmissionHook::evaluate}, which loads the referenced
artifacts, checks the treaty and continuation, verifies the bilateral
statement, and returns either a named denial or an accepted result. Budget
admission runs after the hook returns, so a treaty denial never charges the
capability budget. The tool server is invoked only after these stages succeed.

\subsection{Receiver-Owned Resolution}

Trusted runtime state classifies each request: the receiving edge sets
\codepath{federated_origin_kernel_id}, and the caller cannot. Local requests
follow the local admission path. A request classified as federated must carry
a complete set of identifiers and bindings; malformed or missing context
denies. \codepath{treaty_ref_from_request} reads only identifiers and digests
from the treaty context and resolves every artifact through the hook's own
\codepath{RuntimeAdmissionStores}. That identifier-only parsing plus store
resolution is the structural guarantee that request metadata cannot install
trust material. As an explicit second layer, the parser rejects eleven
trust-shaped keys (trust roots and bundles, treaty scopes, ladder manifests,
signing keys, peer directories, and dynamic trust inputs) with named
request-smuggling codes.

The receiver loads the treaty scope, ladder intersection, continuation,
lineage bundle, bilateral invocation, bilateral DSSE envelope, and the
admission bundle that names the capability lease and governance receipt. It
checks schemas, hashes, validity intervals, participant uniqueness,
action-class coverage, and the receiver's minimum ladder mode. The hook
consumes the continuation during admission evaluation with a single store
operation, records the consumed identifier in the accepted result, and
releases it only on pre-dispatch denial or abort. If release fails, the kernel
keeps the continuation consumed rather than freeing it and records the failure
in the denial receipt.

\subsection{Bilateral Statement}

Three functions verify the bilateral statement before dispatch.
\codepath{verify_treaty_dsse_evidence}, inside the hook, decodes the
statement, requires the strict predicate type
\codepath{chio.bilateral-cosign-invocation.v1}, validates the
policy-evaluation summary and requires an allow verdict, compares the treaty
binding reference (treaty identifier, treaty-scope and ladder-intersection
digests, continuation digest, action class, consistency model, and request
digest) with the artifacts it resolved, checks the tool-argument hash against
that request digest, matches the lease and governance references against its
admission bundle, and requires the signer set to equal the two treaty
participants with distinct public keys. It then calls
\codepath{verify_chio_bilateral_dsse_envelope} and afterwards checks the
lineage-bundle and bilateral-invocation bindings.

\codepath{verify_chio_bilateral_dsse_envelope} enforces the strict DSSE
profile. It checks the payload type, exactly two signatures, byte equality
between the payload and its canonical form, the statement and predicate types,
the predicate schema, one subject carrying the canonical receipt name, two
distinct configured keys with distinct key identifiers and matching
fingerprints, and finally both Ed25519 signatures over the same
pre-authentication encoding. Its failures carry
\codepath{BilateralCoSigningError::code} values such as
\codepath{dsse.malformed}, \codepath{statement.malformed},
\codepath{predicate.type_unrecognised}, and
\codepath{signer.independence_required}.

After \codepath{evaluate_cross_boundary_admission} accepts,
\codepath{VerifiedFederationTreatyMaterial::verify} verifies the envelope
again against the participant keys, checks the participants and tool name,
recomputes the canonical request hash, requires the policy allow result,
checks lease expiry against the injected clock, and compares the treaty,
ladder, action, consistency, and co-sign fields with the accepted admission
report. The cryptographic result is authorization by two configured keys.
Whether those keys have independent organizational custody is outside the
verifier.

\subsection{Binding Admission to the Signed Receipt}

An accepted admission result is converted into
\codepath{VerifiedFederationTreatyMaterial}. Its constructor checks the
request's origin and local participant, distinct participant keys, the DSSE
predicate, the canonical request hash, the policy allow result, the capability
lease, the treaty and ladder hashes, and the action and consistency classes.
It then binds the local accepted admission-report digest into the treaty
binding reference, replacing any peer-supplied value rather than comparing
against it, so a locally signed receipt never carries a peer-supplied report
digest. The kernel stores this object for the current request rather than
retaining a loose collection of request metadata.

Receipt construction rechecks the request, participants, tool, action hash,
and receipt contents against that object. Only then does it attach the treaty
metadata and the resulting outcome and receipt digests to the signed receipt.
This second binding prevents an accepted report from being reused to assemble
a receipt for a different request or participant snapshot.

\subsection{Bounded Predicate Library}

The versioned document type
\codepath{chio.federation.bounded-treaty-predicate.v1} rejects unknown
versions, tags, and fields. It enforces a 64~KiB document limit, a 1~KiB atom
string limit, and depth and node limits of 32 and 1,024.
\codepath{bounded_treaty_receipt_view_from_verified_artifacts} constructs the
input described in Section~\ref{sec:model}. It requires an independently
authenticated admission-report digest and cross-checks the treaty, action,
invocation, and continuation before returning the view.

This library exists for explicit receipt-policy evaluation and differential
testing. It does not grant access on behalf of the admission hook. The hook's
allow decision still depends on the full set of checks above.

\subsection{Runtime and Offline Review}

\codepath{run_runtime_loopback_scenario} drives a deterministic three-vendor
workflow through SQLite receipt and authority stores with a JSON runtime
admission store. It produces the treaty scope, ladder intersection,
continuation, lineage, bilateral DSSE statement, local and remote receipts,
and a buyer package. Producer mode stops after generation and schema
validation. Full mode also invokes the buyer CLI and performs semantic package
review. The offline \codepath{verify_package} path checks the linked
artifacts, hashes, and signatures without access to the vendor processes.

\codepath{verify_package} calls \codepath{verify_chio_bilateral_invocation},
the strict operational verifier, and that is where the remaining
receiver-state comparisons run: peer pins and key fingerprints, the strict
envelope verification, ladder-manifest freshness for both participants, the
revocation epoch, the resolved receipt's signature, kernel key, tool name, and
hashes, the capability lease, policy-verdict agreement, and required
governance evidence. Its failures carry \codepath{VerifierError::code}
values such as \codepath{peer.unpinned_or_keyid_mismatch},
\codepath{peer.revoked_at_epoch}, \codepath{ladder.manifest_stale},
\codepath{capability.lease_expired_or_unknown}, and
\codepath{governance.receipt_required_missing}. The pre-dispatch hook does not
call this verifier.

The negative matrix reaches five surfaces. Five cases (04, 15, 16, 17, and
20) run through \codepath{ChioRuntimeAdmissionHook::evaluate}; three (05,
06, and 07) through the strict operational verifier
\codepath{verify_chio_bilateral_invocation}; six (08 to 12 and 18) through
the envelope verifier or signer, of which only 09 reaches
\codepath{verify_chio_bilateral_dsse_envelope} while the others exercise the
signature-slice envelope verifier or the signer; three (13, 14, and 19)
through the partial verifier \codepath{verify_bilateral_cosign_invocation},
which its own documentation scopes as not fully conformant; and three (01,
02, and 03) through pure runtime-core validators,
\codepath{bounded_treaty_receipt_view_from_verified_artifacts} and
\codepath{evaluate_cross_boundary_admission}, without any DSSE verifier. A
separate dispatch-capable benchmark installs a test tool server and checks
that its invocation counter remains zero after a denial and equals the
iteration count after an admission. Table~\ref{tab:artifacts} lists the
surfaces the evaluation exercises.

\begin{table*}[t]
\caption{Implementation surfaces exercised by the evaluation.}
\label{tab:artifacts}
\scriptsize
\begin{tabularx}{\textwidth}{@{}>{\raggedright\arraybackslash}p{.27\textwidth}>{\raggedright\arraybackslash}p{.37\textwidth}>{\raggedright\arraybackslash}X@{}}
\toprule
Surface & Purpose & Evaluation \\
\midrule
\codepath{ChioRuntimeAdmissionHook::evaluate}
& Receiver-owned resolution, treaty binding checks, and pre-dispatch decision
& Runtime tests and dispatch-capable admission benchmarks \\
\codepath{verify_chio_bilateral_dsse_envelope}
& Strict envelope verification: canonical bytes, two distinct configured keys, both signatures
& Hook and kernel binding paths, federation tests, negative matrix \\
\codepath{verify_chio_bilateral_invocation}
& Offline strict verification with peer pins, manifest freshness, revocation epoch, and receipt store
& Buyer-package verification, federation tests, negative matrix \\
\codepath{VerifiedFederationTreatyMaterial}
& Couples accepted treaty evidence to receipt construction
& Kernel federation and receipt tests \\
\codepath{CrossKernelContinuation}
& Binds parent receipt, nonce, action, target, and lifetime
& Freshness and replay cases \\
\codepath{bounded_treaty_receipt_view_from_verified_artifacts}
& Builds the bounded evaluator input from authenticated artifacts
& Constructor tests and differential suite \\
\codepath{run_runtime_loopback_scenario}
& Produces the complete local buyer package
& End-to-end workflow samples \\
\codepath{verify_package}
& Rechecks the emitted package offline
& Buyer-review tests and component benchmark \\
\bottomrule
\end{tabularx}
\end{table*}

The supplementary manifest records the exact source files, test commands,
benchmark outputs, theorem declarations, and toolchains used for this paper.
It identifies each surface by file and symbol and verifies the content hashes
when the package is rebuilt.
